\section{Data Cleaning and Preprocessing}

Cleaning addressed four concerns: the missing target, predictor imputation,
outliers, and encoding.

\textbf{The missing target is not missing at random.} 21.4\% of Bagrut grade
values are absent, but the absence is structured: suppressed cells have a median
of 7 test takers against 28 where the grade is present. This is the Ministry
protecting the privacy of small cohorts, not random loss. We therefore
\emph{never impute the target}. School-level grades are aggregated over observed
cells only, so no outcome value is invented.

\textbf{Imputation is measured, not assumed.} Rather than assuming multivariate
imputation by chained equations (MICE) recovers missing predictor values well, we
test it directly: we hide 8\% of a fully observed feature (the CBS socioeconomic
score), reconstruct it with MICE, and score the reconstruction against the
withheld truth over 25 independent masks. MICE beats a median baseline on all 25 runs, but the
margin depends on what else is available. With the municipal \texttt{cluster} in
the predictor set MICE reaches $R^2 = 0.954 \pm 0.006$, a number that flatters
the method because \texttt{cluster} is a discretisation of the same index
($r = 0.97$). Withholding that near-duplicate and the outcome columns, MICE still
reaches $R^2 = 0.331 \pm 0.056$ (RMSE $0.755$) while the median baseline never
rises above zero, and it wins all 25 masks in every predictor set tested. MICE
reconstructs each gap from the other observed columns, whereas the median returns
one value everywhere and cannot track variation at all. Per-variant results are
in the repository.

\textbf{Modelling itself uses complete cases.} The masking study is a diagnostic.
Imputed values are \emph{not} carried into the modelling table. Rows missing any
numeric predictor are dropped, costing 9.0\% to 12.5\% of the target-present rows
depending on the outcome, leaving 2,963 to 3,226 school-years per target. Missing
categorical values enter as an all-zero indicator pattern.

\textbf{Outliers are flagged, and kept in the primary analysis.} Two
complementary detectors, Isolation Forest (global anomalies) and the Local
Outlier Factor (local-density anomalies), were run on a nine-feature space. They
agree on only 49 records (Jaccard $0.17$), confirming they capture genuinely
different notions of ``anomalous''. That feature space includes outcome
variables, so excluding the flagged rows would condition the evaluation sample on
the target and could quietly remove legitimate high or low performers. The
primary results therefore retain every valid school. Repeating the evaluation
with the 49 records dropped changes mean $R^2$ by $+0.001$, with per-target
shifts between $-0.023$ and $+0.013$, well inside the fold-to-fold standard
deviation. The choice is immaterial to the findings.

\textbf{Encoding and scaling.} Categorical attributes (sector, supervision,
district, legal status, education stage, locality type) are one-hot encoded, with
Hebrew values mapped to English first for readable outputs. The exam
\texttt{year} is encoded as four independent periods rather than a numeric trend,
since with only four years there is no basis to assume an equal step between
them. Linear models are standardised inside the cross-validation folds, never on
the full dataset, so no test-fold information leaks into the scaling. Trees need
no scaling.
